\begin{lexlean}{Uorc.Core.RuntimeAdmissionCreditBoundary}
\importmodule{Uorc.Core.RuntimeAdmissionCredit}
\importmodule{Uorc.Core.RuntimeAdmission}
\importmodule{Uorc.Core.RuntimeAdmissionBoundary}
\importmodule{Uorc.Core.RuntimeResources}
\importmodule{Uorc.Core.RuntimeProgram}
\importmodule{Uorc.Core.Heap}
\importmodule{Uorc.Core.RuntimeValue}
\title{Parameters}
\begin{semanticmodule}
structure{
  name: Accounting,
  type_parameters: [],
  parameters: [],
  fields: [{name: consumedCredits, type: UInt64}, {name: attemptedOperations, type: UInt64}]
}

structure{
  name: Detached,
  type_parameters: [],
  parameters: [],
  fields: [
    {name: state, type: Uorc.Core.RuntimeAdmissionBoundary::AwaitingState},
    {name: accounting, type: Accounting}
  ]
}

structure{
  name: Taken,
  type_parameters: [],
  parameters: [],
  fields: [
    {name: arena, type: Uorc.Core.Heap::Arena},
    {name: usage, type: Uorc.Core.RuntimeResources::Usage},
    {name: caller, type: Uorc.Core.RuntimeProgram::CallerRoots},
    {name: session, type: Detached}
  ]
}

structure{
  name: Returning,
  type_parameters: [],
  parameters: [],
  fields: [
    {name: boundary, type: Uorc.Core.RuntimeAdmissionBoundary::ReturnCursor},
    {name: accounting, type: Accounting}
  ]
}

inductive{
  name: Error,
  type_parameters: [],
  parameters: [],
  constructors: [
    {name: busy, fields: []},
    {name: boundary, fields: [Uorc.Core.RuntimeAdmissionBoundary::Error]},
    {name: credit, fields: [Uorc.Core.RuntimeAdmissionCredit::Error]}
  ]
}

inductive{
  name: TakeEvent,
  type_parameters: [],
  parameters: [],
  constructors: [
    {name: taken, fields: [Taken]},
    {name: failed, fields: [Error, Uorc.Core.RuntimeAdmissionCredit::Cursor]},
    {name: retainedFailure, fields: [Error, Taken]}
  ]
}

inductive{
  name: ReturnEvent,
  type_parameters: [],
  parameters: [],
  constructors: [
    {name: pending, fields: [Returning]},
    {name: restored, fields: [Uorc.Core.RuntimeAdmissionCredit::Cursor]},
    {name: failed, fields: [Error, Returning]}
  ]
}

definition{
  name: take,
  parameters: [
    cursor : Uorc.Core.RuntimeAdmissionCredit::Cursor,
    role : Uorc.Core.RuntimeAdmissionBoundary::ProducerRole
  ],
  result: TakeEvent,
  body: (match (project cursor error)
    | Option.none => (if ((prim equal (project cursor activeCost) (0 : UInt64) : Bool) && (prim equal (project cursor paidCost) (0 : UInt64) : Bool)) then (match (project cursor intent)
      | Uorc.Core.RuntimeAdmissionCredit::Intent.initialize budget arena basis caller usage => (ctor TakeEvent.failed (ctor Error.busy) cursor)
      | Uorc.Core.RuntimeAdmissionCredit::Intent.execute state => (ctor TakeEvent.failed (ctor Error.busy) cursor)
      | Uorc.Core.RuntimeAdmissionCredit::Intent.offer state event => (ctor TakeEvent.failed (ctor Error.busy) cursor)
      | Uorc.Core.RuntimeAdmissionCredit::Intent.terminal outcome => (match outcome
        | Uorc.Core.RuntimeAdmission::Step.pending state => (ctor TakeEvent.failed (ctor Error.busy) cursor)
        | Uorc.Core.RuntimeAdmission::Step.awaiting state => (match (Uorc.Core.RuntimeAdmissionBoundary::take state role)
          | Uorc.Core.RuntimeAdmissionBoundary::TakeEvent.taken loan => (ctor TakeEvent.taken (record Taken {arena := (project loan arena), usage := (project loan usage), caller := (project loan caller), session := (record Detached {state := (project loan state), accounting := (record Accounting {consumedCredits := (project cursor consumedCredits), attemptedOperations := (project cursor attemptedOperations)})})}))
          | Uorc.Core.RuntimeAdmissionBoundary::TakeEvent.failed failure next => (ctor TakeEvent.failed (ctor Error.boundary failure) (record Uorc.Core.RuntimeAdmissionCredit::Cursor {intent := (ctor Uorc.Core.RuntimeAdmissionCredit::Intent.terminal (ctor Uorc.Core.RuntimeAdmission::Step.awaiting next)), activeCost := (0 : UInt64), paidCost := (0 : UInt64), consumedCredits := (project (record Accounting {consumedCredits := (project cursor consumedCredits), attemptedOperations := (project cursor attemptedOperations)}) consumedCredits), attemptedOperations := (project (record Accounting {consumedCredits := (project cursor consumedCredits), attemptedOperations := (project cursor attemptedOperations)}) attemptedOperations), usage := (project next usage), error := (ctor Option.none [Uorc.Core.RuntimeAdmissionCredit::Error])}))
          | Uorc.Core.RuntimeAdmissionBoundary::TakeEvent.retainedFailure failure loan => (ctor TakeEvent.retainedFailure (ctor Error.boundary failure) (record Taken {arena := (project loan arena), usage := (project loan usage), caller := (project loan caller), session := (record Detached {state := (project loan state), accounting := (record Accounting {consumedCredits := (project cursor consumedCredits), attemptedOperations := (project cursor attemptedOperations)})})})))
        | Uorc.Core.RuntimeAdmission::Step.failed failure state => (ctor TakeEvent.failed (ctor Error.busy) cursor)
        | Uorc.Core.RuntimeAdmission::Step.ready ready => (ctor TakeEvent.failed (ctor Error.busy) cursor))) else (ctor TakeEvent.failed (ctor Error.busy) cursor))
    | Option.some failure => (ctor TakeEvent.failed (ctor Error.credit failure) cursor)),
  axioms: [Quot.sound, propext],
  executable: true
}

definition{
  name: reattach,
  parameters: [
    session : Detached,
    role : Uorc.Core.RuntimeAdmissionBoundary::ProducerRole,
    arena : Uorc.Core.Heap::Arena,
    usage : Uorc.Core.RuntimeResources::Usage,
    caller : Uorc.Core.RuntimeProgram::CallerRoots
  ],
  result: ReturnEvent,
  body: (match (Uorc.Core.RuntimeAdmissionBoundary::reattach (project session state) role arena usage caller)
    | Uorc.Core.RuntimeAdmissionBoundary::ReturnEvent.pending pending => (ctor ReturnEvent.pending (record Returning {boundary := pending, accounting := (project session accounting)}))
    | Uorc.Core.RuntimeAdmissionBoundary::ReturnEvent.restored state => (ctor ReturnEvent.restored (record Uorc.Core.RuntimeAdmissionCredit::Cursor {intent := (ctor Uorc.Core.RuntimeAdmissionCredit::Intent.terminal (ctor Uorc.Core.RuntimeAdmission::Step.awaiting state)), activeCost := (0 : UInt64), paidCost := (0 : UInt64), consumedCredits := (project (project session accounting) consumedCredits), attemptedOperations := (project (project session accounting) attemptedOperations), usage := (project state usage), error := (ctor Option.none [Uorc.Core.RuntimeAdmissionCredit::Error])}))
    | Uorc.Core.RuntimeAdmissionBoundary::ReturnEvent.failed failure pending => (ctor ReturnEvent.failed (ctor Error.boundary failure) (record Returning {boundary := pending, accounting := (project session accounting)}))),
  axioms: [Quot.sound, propext],
  executable: true
}

definition{
  name: returnStep,
  parameters: [returning : Returning],
  result: ReturnEvent,
  body: (match (Uorc.Core.RuntimeAdmissionBoundary::returnStep (project returning boundary))
    | Uorc.Core.RuntimeAdmissionBoundary::ReturnEvent.pending pending => (ctor ReturnEvent.pending (record Returning {boundary := pending, accounting := (project returning accounting)}))
    | Uorc.Core.RuntimeAdmissionBoundary::ReturnEvent.restored state => (ctor ReturnEvent.restored (record Uorc.Core.RuntimeAdmissionCredit::Cursor {intent := (ctor Uorc.Core.RuntimeAdmissionCredit::Intent.terminal (ctor Uorc.Core.RuntimeAdmission::Step.awaiting state)), activeCost := (0 : UInt64), paidCost := (0 : UInt64), consumedCredits := (project (project returning accounting) consumedCredits), attemptedOperations := (project (project returning accounting) attemptedOperations), usage := (project state usage), error := (ctor Option.none [Uorc.Core.RuntimeAdmissionCredit::Error])}))
    | Uorc.Core.RuntimeAdmissionBoundary::ReturnEvent.failed failure pending => (ctor ReturnEvent.failed (ctor Error.boundary failure) (record Returning {boundary := pending, accounting := (project returning accounting)}))),
  axioms: [Quot.sound, propext],
  executable: true
}

definition{
  name: takenRootInventory,
  parameters: [taken : Taken],
  result: (List Uorc.Core.RuntimeValue::Handle),
  body: (Uorc.Core.RuntimeAdmissionBoundary::rootInventory (project (project taken session) state) (project taken caller)),
  axioms: [],
  executable: true
}

definition{
  name: detachedRootInventory,
  parameters: [session : Detached, caller : Uorc.Core.RuntimeProgram::CallerRoots],
  result: (List Uorc.Core.RuntimeValue::Handle),
  body: (Uorc.Core.RuntimeAdmissionBoundary::rootInventory (project session state) caller),
  axioms: [],
  executable: true
}

definition{
  name: returnRootInventory,
  parameters: [returning : Returning],
  result: (List Uorc.Core.RuntimeValue::Handle),
  body: (Uorc.Core.RuntimeAdmissionBoundary::returnRootInventory (project returning boundary)),
  axioms: [],
  executable: true
}
\end{semanticmodule}
\end{lexlean}
